\documentclass[10pt,a4paper]{amsart}
\usepackage{amsmath,amssymb,amsthm,graphicx}
% Retained minimal preamble for the Stacks Project chapter varieties.tex.
% Verbatim excerpt of the chapter preamble: Xy-pic setup (native preamble lines 14-18), the native theorem
% declarations with the mathematical macros (native preamble lines 154-217), and the project own no-op
% definitions of the reference/slogan/history display environments (native wrapper.tex lines 1-3) that the
% counted statement uses for its slogan. Copied mechanically; nothing added, removed or redefined.
% For commutative diagrams we use Xy-pic
\usepackage[all]{xy}

% We use 2cell for 2-commutative diagrams.
\xyoption{2cell}
% Theorem environments.
%
\theoremstyle{plain}
\newtheorem{theorem}[subsection]{Theorem}
\newtheorem{proposition}[subsection]{Proposition}
\newtheorem{lemma}[subsection]{Lemma}

\theoremstyle{definition}
\newtheorem{definition}[subsection]{Definition}
\newtheorem{example}[subsection]{Example}
\newtheorem{exercise}[subsection]{Exercise}
\newtheorem{situation}[subsection]{Situation}

\theoremstyle{remark}
\newtheorem{remark}[subsection]{Remark}
\newtheorem{remarks}[subsection]{Remarks}

\numberwithin{equation}{subsection}

% Macros
%
\def\lim{\mathop{\mathrm{lim}}\nolimits}
\def\colim{\mathop{\mathrm{colim}}\nolimits}
\def\Spec{\mathop{\mathrm{Spec}}}
\def\Hom{\mathop{\mathrm{Hom}}\nolimits}
\def\Ext{\mathop{\mathrm{Ext}}\nolimits}
\def\SheafHom{\mathop{\mathcal{H}\!\mathit{om}}\nolimits}
\def\SheafExt{\mathop{\mathcal{E}\!\mathit{xt}}\nolimits}
\def\Sch{\mathit{Sch}}
\def\Mor{\mathop{\mathrm{Mor}}\nolimits}
\def\Ob{\mathop{\mathrm{Ob}}\nolimits}
\def\Sh{\mathop{\mathit{Sh}}\nolimits}
\def\NL{\mathop{N\!L}\nolimits}
\def\CH{\mathop{\mathrm{CH}}\nolimits}
\def\proetale{{pro\text{-}\acute{e}tale}}
\def\etale{{\acute{e}tale}}
\def\QCoh{\mathit{QCoh}}
\def\Ker{\mathop{\mathrm{Ker}}}
\def\Im{\mathop{\mathrm{Im}}}
\def\Coker{\mathop{\mathrm{Coker}}}
\def\Coim{\mathop{\mathrm{Coim}}}

% Boxtimes
%
\DeclareMathSymbol{\boxtimes}{\mathbin}{AMSa}{"02}

%
% Macros for moduli stacks/spaces
%
\def\QCohstack{\mathcal{QC}\!\mathit{oh}}
\def\Cohstack{\mathcal{C}\!\mathit{oh}}
\def\Spacesstack{\mathcal{S}\!\mathit{paces}}
\def\Quotfunctor{\mathrm{Quot}}
\def\Hilbfunctor{\mathrm{Hilb}}
\def\Curvesstack{\mathcal{C}\!\mathit{urves}}
\def\Polarizedstack{\mathcal{P}\!\mathit{olarized}}
\def\Complexesstack{\mathcal{C}\!\mathit{omplexes}}
% \Pic is the operator that assigns to X its picard group, usage \Pic(X)
% \Picardstack_{X/B} denotes the Picard stack of X over B
% \Picardfunctor_{X/B} denotes the Picard functor of X over B
\def\Pic{\mathop{\mathrm{Pic}}\nolimits}
\def\Picardstack{\mathcal{P}\!\mathit{ic}}
\def\Picardfunctor{\mathrm{Pic}}
\def\Deformationcategory{\mathcal{D}\!\mathit{ef}}
\newenvironment{reference}{}{}
\newenvironment{slogan}{}{}
\newenvironment{history}{}{}
\usepackage{cleveref}
\crefname{originallemma}{original lemma}{original lemmas}
\Crefname{originallemma}{Original lemma}{Original lemmas}
\makeatletter
\newlabel{lemma-dimension-fibre-in-codim-1}{{0B2I}{source}{}{sourcepointer-1}{}}
\newlabel{lemma-dimension-fibre-in-codim-1@cref}{{[originallemma][0][]0B2I}{[1][0][]source}{}{}{}}
\makeatother
\crefname{originallemma}{original lemma}{original lemmas}
\Crefname{originallemma}{Original lemma}{Original lemmas}
\makeatletter
\newlabel{lemma-dimension-fibre-in-higher-codimension}{{0B2K}{source}{}{sourcepointer-2}{}}
\newlabel{lemma-dimension-fibre-in-higher-codimension@cref}{{[originallemma][0][]0B2K}{[1][0][]source}{}{}{}}
\makeatother
\crefname{originallemma}{original lemma}{original lemmas}
\Crefname{originallemma}{Original lemma}{Original lemmas}
\makeatletter
\newlabel{lemma-dimension-locally-algebraic}{{0A21}{source}{}{sourcepointer-3}{}}
\newlabel{lemma-dimension-locally-algebraic@cref}{{[originallemma][0][]0A21}{[1][0][]source}{}{}{}}
\makeatother
\crefname{originallemma}{original lemma}{original lemmas}
\Crefname{originallemma}{Original lemma}{Original lemmas}
\makeatletter
\newlabel{morphisms-lemma-openness-bounded-dimension-fibres}{{02FZ}{source}{}{sourcepointer-4}{}}
\newlabel{morphisms-lemma-openness-bounded-dimension-fibres@cref}{{[originallemma][0][]02FZ}{[1][0][]source}{}{}{}}
\makeatother
\crefname{originallemma}{original lemma}{original lemmas}
\Crefname{originallemma}{Original lemma}{Original lemmas}
\makeatletter
\newlabel{morphisms-lemma-permanence-finite-type}{{01T8}{source}{}{sourcepointer-5}{}}
\newlabel{morphisms-lemma-permanence-finite-type@cref}{{[originallemma][0][]01T8}{[1][0][]source}{}{}{}}
\makeatother
\title{Dimension of fibres of a morphism of locally algebraic schemes: bounds, openness and the generic value}
\author{The Stacks Project Authors}
\date{stacks-project:varieties.tex:local-snapshot-20261011; GNU FDL 1.2 or later}
\begin{document}
\maketitle
\noindent\textit{Editorial note.} All delivered mathematical text is the author's own, copied from the pinned Stacks Project chapter file varieties.tex; no statement or proof step was rewritten, reordered or removed.\par\medskip
\noindent\textit{Editorial note.} The proof states that the author will freely use the results of the dimension theory lemma for X, Y and the fibres; the sentence is reproduced verbatim and those results are navigation stubs, while the computation of d is written out inside the counted span.\par\medskip
\noindent\textit{Editorial note.} The counted spans carry 6 navigation stubs and no bibliography citation.\par\medskip
\noindent\textit{Editorial note.} The preamble is the chapter's retained minimal preamble (sources/stacks-varieties/retained-preamble.tex lines 1-77), a verbatim excerpt of the native preamble.\par\medskip
\noindent\textit{Editorial note.} The label lemma-dimension-fibres-locally-algebraic is defined only in varieties.tex, so no cross-chapter restatement had to be weighed for this card.\par\medskip
\noindent\textit{Allowed standard prerequisites (editorial).}\begin{itemize}
\item Transcendence degree description of the dimension of a scheme and of its fibres. Same chapter, Lemma lemma-dimension-locally-algebraic, used freely at proof lines 3979-3981 and in the computation at 3986-3990. Supplies the identification of d with the fibre dimension at the generic point.
\item Openness of the locus where the fibre dimension is bounded. Morphisms, Lemma morphisms-lemma-openness-bounded-dimension-fibres, invoked at proof lines 3991-3992. Supplies the lower bound and the density of equality.
\item Fibre dimension bounds in codimension one and in higher codimension. Same chapter, Lemmas lemma-dimension-fibre-in-codim-1 and lemma-dimension-fibre-in-higher-codimension, invoked at proof lines 3993-3995. Supplies parts (2)(c) and (2)(d).
\item Locally finite type is preserved by the morphism. Morphisms, Lemma morphisms-lemma-permanence-finite-type, invoked at proof lines 3975-3978. Supplies that the fibres are locally algebraic over the residue field.
\end{itemize}\medskip
\ifdefined\hypertarget\hypertarget{sourcepointer-1}{}\fi
\noindent\textit{Source reference.} Original lemma 0B2I: sources/stacks-varieties/pinned-original-source.tex, source lines 3597--3597. Pointer outside the counted span, invoked once for context; the load-bearing argument of the counted statement is the author's proof reproduced in full in the proof span.\par\medskip
\ifdefined\hypertarget\hypertarget{sourcepointer-2}{}\fi
\noindent\textit{Source reference.} Original lemma 0B2K: sources/stacks-varieties/pinned-original-source.tex, source lines 3671--3671. Pointer outside the counted span, invoked once for context; the load-bearing argument of the counted statement is the author's proof reproduced in full in the proof span.\par\medskip
\ifdefined\hypertarget\hypertarget{sourcepointer-3}{}\fi
\noindent\textit{Source reference.} Original lemma 0A21: sources/stacks-varieties/pinned-original-source.tex, source lines 3783--3783. Pointer outside the counted span, invoked once for context; the load-bearing argument of the counted statement is the author's proof reproduced in full in the proof span.\par\medskip
\ifdefined\hypertarget\hypertarget{sourcepointer-4}{}\fi
\noindent\textit{Source reference.} Original lemma 02FZ: sources/stacks-chow/tags.txt, source lines 3210--3210. Pointer outside the counted span, invoked once for context; the load-bearing argument of the counted statement is the author's proof reproduced in full in the proof span.\par\medskip
\ifdefined\hypertarget\hypertarget{sourcepointer-5}{}\fi
\noindent\textit{Source reference.} Original lemma 01T8: sources/stacks-chow/tags.txt, source lines 2365--2365. Pointer outside the counted span, invoked once for context; the load-bearing argument of the counted statement is the author's proof reproduced in full in the proof span.\par\medskip
\begin{lemma}
\label{lemma-dimension-fibres-locally-algebraic}
Let $k$ be a field. Let $f : X \to Y$ be a morphism of locally algebraic
$k$-schemes.
\begin{enumerate}
\item For $y \in Y$, the fibre $X_y$ is a locally
algebraic scheme over $\kappa(y)$ hence all the results of
Lemma \ref{lemma-dimension-locally-algebraic} apply.
\item Assume $X$ is irreducible. Set $Z = \overline{f(X)}$ and
$d = \dim(X) - \dim(Z)$. Then
\begin{enumerate}
\item $\dim_x(X_{f(x)}) \geq d$ for all $x \in X$,
\item the set of $x \in X$ with $\dim_x(X_{f(x)}) = d$ is dense open,
\item if $\dim(\mathcal{O}_{Z, f(x)}) \geq 1$, then
$\dim_x(X_{f(x)}) \leq d + \dim(\mathcal{O}_{Z, f(x)}) - 1$,
\item if $\dim(\mathcal{O}_{Z, f(x)}) = 1$, then $\dim_x(X_{f(x)}) = d$,
\end{enumerate}
\item For $x \in X$ with $y = f(x)$ we have
$\dim_x(X_y) \geq \dim_x(X) - \dim_y(Y)$.
\end{enumerate}
\end{lemma}

\begin{proof}
The morphism $f$ is locally of finite type by
Morphisms, Lemma \ref{morphisms-lemma-permanence-finite-type}.
Hence the base change $X_y \to \Spec(\kappa(y))$ is locally
of finite type. This proves (1).
In the rest of the proof we will freely use the results
of Lemma \ref{lemma-dimension-locally-algebraic} for $X$, $Y$, and
the fibres of $f$.

\medskip\noindent
Proof of (2). Let $\eta \in X$ be the generic point and set
$\xi = f(\eta)$. Then $Z = \overline{\{\xi\}}$. Hence
$$
d = \dim(X) - \dim(Z) =
\text{trdeg}_k \kappa(\eta) - \text{trdeg}_k \kappa(\xi) =
\text{trdeg}_{\kappa(\xi)} \kappa(\eta) = \dim_\eta(X_\xi)
$$
Thus parts (2)(a) and (2)(b) follow from
Morphisms, Lemma \ref{morphisms-lemma-openness-bounded-dimension-fibres}.
Parts (2)(c) and (2)(d) follow from
Lemmas \ref{lemma-dimension-fibre-in-higher-codimension} and
\ref{lemma-dimension-fibre-in-codim-1}.

\medskip\noindent
Proof of (3). Let $x \in X$. Let $X' \subset X$ be an irreducible component
of $X$ passing through $x$ of dimension $\dim_x(X)$. Then (2) implies that
$\dim_x(X_y) \geq \dim(X') - \dim(Z')$ where $Z' \subset Y$
is the closure of the image of $X'$. This proves (3).
\end{proof}

\end{document}
